\begin{table}[h]
\centering
\caption{Characterization of knowledge graphs and advantages in representational learning \parencite{yu2024review}.}
\begin{tabular}{|p{0.45\textwidth}|p{0.45\textwidth}|}
\hline
\textbf{Characteristics of the Knowledge Graph} & \textbf{Indicates strengths in recommendation} \\
\hline
1) Structured representation: representation of knowledge in the form of diagrams, which can clearly express the relationship between entities and facilitate understanding and processing. & 1) Improve recommendation accuracy and personalization: using entity and relationship information from the knowledge graph, recommender systems can more accurately understand the relationship between users and items, thus improving recommendation accuracy and personalization. \\
\hline
2) Semantically rich: contains not only entities and relationships but also attribute information of entities, able to express rich semantic information. & 2) Solving the cold-start problem: The knowledge graph solves the cold-start problem by introducing attribute information about users and projects so that the recommender system can better adapt to new users and new projects. \\
\hline
3) Extensibility: good extensibility to easily add new entities, relationships, and attributes to accommodate growing knowledge and information. & 3) Complex relationship modeling: knowledge graphs can represent complex relationships and attribute information, which helps recommender systems capture the interaction patterns between users and items more accurately and improves recommendation accuracy and personalization. \\
\hline
4) Dynamic updates: can be dynamically updated and maintained to reflect real-world changes and updates. & \\
\hline
\end{tabular}
\end{table}